% =====================================================================
%  BUỔI 1 — TÀI LIỆU BUỔI HỌC (từ 02-tai-lieu-buoi-hoc.md)
%  Học sinh dùng trên lớp và tra cứu khi làm bài.
% =====================================================================
\documentclass[11pt, a4paper]{article}
% =====================================================================
%  CAE SHANGHAI — GÓI LỆNH DÙNG CHUNG CHO TÀI LIỆU LUYỆN THI CSCA TOÁN
%  Biên dịch bằng XeLaTeX. Mọi file .tex trong buổi học \input file này.
% =====================================================================

% ---------- NGÔN NGỮ & PHÔNG CHỮ ----------
\usepackage{fontspec}
\usepackage{polyglossia}
\setmainlanguage{vietnamese}
\setotherlanguage{english}

% Chữ Latin (Anh + Việt): TeX Gyre Termes — bản clone Times New Roman,
% đi kèm TeX Live, hỗ trợ đầy đủ dấu tiếng Việt.
\setmainfont{TeX Gyre Termes}
\setsansfont{TeX Gyre Heros}

% Chữ Hán: ưu tiên Source Han Serif SC (思源宋体) — chuẩn học thuật Trung Quốc.
% Nếu máy chưa cài, tự động dùng Songti SC (宋体, có sẵn trên macOS).
% Cả hai đều thuộc họ 宋体 nên diện mạo tài liệu không đổi.
\IfFontExistsTF{Source Han Serif SC}{
  \newcommand{\cjkfamilyname}{Source Han Serif SC}
}{
  \IfFontExistsTF{Songti SC}{
    \newcommand{\cjkfamilyname}{Songti SC}
  }{
    \newcommand{\cjkfamilyname}{Noto Serif CJK SC}
  }
}

% xeCJK tự chuyển font khi gặp ký tự Hán, không cần bọc \cjk{} thủ công.
% Đây là điều kiện bắt buộc: font Latin (TeX Gyre Termes) không có glyph chữ Hán,
% nếu không khai báo xeCJK thì mọi chữ Hán viết trần sẽ bị nuốt khỏi bản in.
\usepackage{xeCJK}
\setCJKfamilyfont{caehan}{\cjkfamilyname}[Scale=0.95]
\setCJKmainfont{\cjkfamilyname}[Scale=0.95]
\setCJKsansfont{\cjkfamilyname}[Scale=0.95]
\setCJKmonofont{\cjkfamilyname}[Scale=0.95]

% \cjk{} ép dùng font Hán trong ngữ cảnh cần thiết (ví dụ trong tiêu đề in đậm
% hoặc trong bảng có định dạng riêng). Dùng \CJKfamily chứ không phải tên font,
% vì tên font chỉ là văn bản — sẽ in ra chữ thay vì đổi font.
\newcommand{\cjk}[1]{{\CJKfamily{caehan}#1}}

% ---------- TOÁN HỌC ----------
\usepackage{amsmath, amssymb, amsthm, mathtools}
\usepackage{unicode-math}
\setmathfont{TeX Gyre Termes Math}

% ---------- BỐ CỤC TRANG ----------
% Lề phải chừa vùng an toàn của letterhead:
%   banner trên chiếm 0–19mm  → top tối thiểu 2,2cm
%   banner dưới chiếm 281–297mm → bottom tối thiểu 2,2cm
% Lề trên 2,6cm (26mm) để tiêu đề đầu trang không đè banner (banner hết ở 19mm).
\usepackage[a4paper, top=2.6cm, bottom=2.4cm, left=2.5cm, right=2.0cm]{geometry}
\usepackage{fancyhdr}
\usepackage{titlesec}
\usepackage{enumitem}
\usepackage{graphicx}
\usepackage{booktabs, array, longtable, tabularx}
\usepackage{xcolor}
\usepackage[most]{tcolorbox}
\usepackage[colorlinks=true, linkcolor=cae_blue, citecolor=cae_blue, urlcolor=cae_blue]{hyperref}

% ---------- BẢNG MÀU THƯƠNG HIỆU ----------
\definecolor{cae_blue}{RGB}{0, 51, 102}
\definecolor{cae_gold}{RGB}{196, 149, 60}
\definecolor{box_bg}{RGB}{245, 247, 250}
\definecolor{box_line}{RGB}{0, 51, 102}

% ---------- TIÊU ĐỀ MỤC ----------
% Đệm đầu trang: banner letterhead chiếm 0–19mm, mà lề trên là 22mm.
% \section* đầu trang (sau \clearpage) sát lề nên phải chừa thêm, nếu không
% sẽ đè lên banner.
\titleformat{\section}
  {\normalfont\large\bfseries\sffamily\color{cae_blue}}{\thesection.}{0.6em}{}
\titleformat{\subsection}
  {\normalfont\normalsize\bfseries\sffamily\color{cae_blue!85!black}}{\thesubsection}{0.6em}{}
\titlespacing*{\section}{0pt}{14pt}{6pt}
\titlespacing*{\subsection}{0pt}{10pt}{4pt}

% Tiêu đề không đánh số ở đầu trang: \section* sau \clearpage nằm sát lề trên
% nên đè lên banner letterhead. Dùng \vspace* ngay TRƯỚC \section* không đủ,
% vì titlesec chèn khoảng cách riêng. Cách chắc chắn là đệm bằng \null.
\newcommand{\sectiondautrang}[1]{\section*{#1}}

% Dạng bài dùng \subsection*{} nên cần định dạng riêng, đậm và có gạch chân màu
\titleformat{\subsubsection}
  {\normalfont\normalsize\bfseries\sffamily\color{cae_blue!85!black}}{\thesubsubsection}{0.6em}{}

% ---------- LETTERHEAD (ảnh nền toàn trang) ----------
% Ảnh letterhead của CAE Shanghai: banner xanh trên/dưới kèm logo chìm giữa.
% Tỉ lệ ảnh 1414x2000 = 0.7070, khớp A4 (0.7071) nên chèn thẳng không méo.
% Đặt ở lớp nền dưới cùng để chữ không bị ảnh che.
% Dùng \AddToShipoutPictureBG (không có dấu *) để áp dụng cho MỌI trang.
\usepackage{eso-pic}
\AddToShipoutPictureBG{%
  \AtPageLowerLeft{%
    \includegraphics[width=\paperwidth,height=\paperheight]{images/letterhead.png}%
  }%
}

% ---------- ĐẦU / CHÂN TRANG ----------
% Letterhead đã có banner trên và dưới, nên bỏ dòng chạy để tránh chồng chữ.
% Số trang đặt trong vùng nội dung, sát trên banner dưới, lề phải.
\pagestyle{fancy}
\setlength{\headheight}{0pt}
\fancyhf{}
\renewcommand{\headrulewidth}{0pt}
\renewcommand{\footrulewidth}{0pt}
\fancyfoot[R]{\raisebox{0.9cm}{\small\color{cae_blue}\thepage}}

% ---------- THỤT LỀ ĐOẠN ----------
% Mặc định LaTeX thụt lề 15pt nhưng bỏ qua đoạn ngay sau tiêu đề mục, nên
% trong cùng một mục có đoạn thụt đoạn không, nhìn lệch hàng. Tài liệu này
% ngắt đoạn bằng khoảng trắng chứ không bằng thụt lề, nên đặt parindent = 0
% và cho parskip giãn ra để vẫn phân biệt được các đoạn.
\setlength{\parindent}{0pt}
\setlength{\parskip}{0.6em}

% ---------- MÔI TRƯỜNG NỘI DUNG ----------
% Trong tcolorbox, LaTeX đặt parindent = parskip = 0, nên các đoạn văn trong
% khung dính liền nhau. Đặt lại khoảng cách đoạn cho mọi khung.
\newcommand{\khoangcachdoan}{\setlength{\parskip}{0.5em}\setlength{\parindent}{0pt}}
\tcbset{before upper={\khoangcachdoan}}

\newtcolorbox{dinhnghia}[1][]{
  colback=box_bg, colframe=box_line, boxrule=0.6pt, arc=1pt,
  fonttitle=\bfseries, title={Định nghĩa #1}, breakable}
\newtcolorbox{dinhly}[1][]{
  colback=box_bg, colframe=cae_gold, boxrule=0.6pt, arc=1pt,
  fonttitle=\bfseries, title={Tính chất #1}, breakable}
\newtcolorbox{vidu}[1][]{
  colback=white, colframe=black!40, boxrule=0.4pt, arc=1pt,
  fonttitle=\bfseries, title={Ví dụ #1}, breakable}
\newtcolorbox{baitap}[1][]{
  colback=white, colframe=cae_blue!50, boxrule=0.4pt, arc=1pt,
  fonttitle=\bfseries, title={Câu #1}, breakable}

% Ô đáp án: viền xanh đậm hơn, dùng trong file đáp án
\newtcolorbox{dapan}[1][]{
  colback=box_bg, colframe=cae_blue, boxrule=0.6pt, arc=1pt,
  fonttitle=\bfseries, title={Câu #1}, breakable}

% ---------- TÊN CHỨNG MINH ----------
\renewcommand{\proofname}{\normalfont\bfseries Chứng minh}

% ---------- MÔN HỌC ----------
% Đổi được ở file .tex trước khi % =====================================================================
%  CAE SHANGHAI — GÓI LỆNH DÙNG CHUNG CHO TÀI LIỆU LUYỆN THI CSCA TOÁN
%  Biên dịch bằng XeLaTeX. Mọi file .tex trong buổi học \input file này.
% =====================================================================

% ---------- NGÔN NGỮ & PHÔNG CHỮ ----------
\usepackage{fontspec}
\usepackage{polyglossia}
\setmainlanguage{vietnamese}
\setotherlanguage{english}

% Chữ Latin (Anh + Việt): TeX Gyre Termes — bản clone Times New Roman,
% đi kèm TeX Live, hỗ trợ đầy đủ dấu tiếng Việt.
\setmainfont{TeX Gyre Termes}
\setsansfont{TeX Gyre Heros}

% Chữ Hán: ưu tiên Source Han Serif SC (思源宋体) — chuẩn học thuật Trung Quốc.
% Nếu máy chưa cài, tự động dùng Songti SC (宋体, có sẵn trên macOS).
% Cả hai đều thuộc họ 宋体 nên diện mạo tài liệu không đổi.
\IfFontExistsTF{Source Han Serif SC}{
  \newcommand{\cjkfamilyname}{Source Han Serif SC}
}{
  \IfFontExistsTF{Songti SC}{
    \newcommand{\cjkfamilyname}{Songti SC}
  }{
    \newcommand{\cjkfamilyname}{Noto Serif CJK SC}
  }
}

% xeCJK tự chuyển font khi gặp ký tự Hán, không cần bọc \cjk{} thủ công.
% Đây là điều kiện bắt buộc: font Latin (TeX Gyre Termes) không có glyph chữ Hán,
% nếu không khai báo xeCJK thì mọi chữ Hán viết trần sẽ bị nuốt khỏi bản in.
\usepackage{xeCJK}
\setCJKfamilyfont{caehan}{\cjkfamilyname}[Scale=0.95]
\setCJKmainfont{\cjkfamilyname}[Scale=0.95]
\setCJKsansfont{\cjkfamilyname}[Scale=0.95]
\setCJKmonofont{\cjkfamilyname}[Scale=0.95]

% \cjk{} ép dùng font Hán trong ngữ cảnh cần thiết (ví dụ trong tiêu đề in đậm
% hoặc trong bảng có định dạng riêng). Dùng \CJKfamily chứ không phải tên font,
% vì tên font chỉ là văn bản — sẽ in ra chữ thay vì đổi font.
\newcommand{\cjk}[1]{{\CJKfamily{caehan}#1}}

% ---------- TOÁN HỌC ----------
\usepackage{amsmath, amssymb, amsthm, mathtools}
\usepackage{unicode-math}
\setmathfont{TeX Gyre Termes Math}

% ---------- BỐ CỤC TRANG ----------
% Lề phải chừa vùng an toàn của letterhead:
%   banner trên chiếm 0–19mm  → top tối thiểu 2,2cm
%   banner dưới chiếm 281–297mm → bottom tối thiểu 2,2cm
% Lề trên 2,6cm (26mm) để tiêu đề đầu trang không đè banner (banner hết ở 19mm).
\usepackage[a4paper, top=2.6cm, bottom=2.4cm, left=2.5cm, right=2.0cm]{geometry}
\usepackage{fancyhdr}
\usepackage{titlesec}
\usepackage{enumitem}
\usepackage{graphicx}
\usepackage{booktabs, array, longtable, tabularx}
\usepackage{xcolor}
\usepackage[most]{tcolorbox}
\usepackage[colorlinks=true, linkcolor=cae_blue, citecolor=cae_blue, urlcolor=cae_blue]{hyperref}

% ---------- BẢNG MÀU THƯƠNG HIỆU ----------
\definecolor{cae_blue}{RGB}{0, 51, 102}
\definecolor{cae_gold}{RGB}{196, 149, 60}
\definecolor{box_bg}{RGB}{245, 247, 250}
\definecolor{box_line}{RGB}{0, 51, 102}

% ---------- TIÊU ĐỀ MỤC ----------
% Đệm đầu trang: banner letterhead chiếm 0–19mm, mà lề trên là 22mm.
% \section* đầu trang (sau \clearpage) sát lề nên phải chừa thêm, nếu không
% sẽ đè lên banner.
\titleformat{\section}
  {\normalfont\large\bfseries\sffamily\color{cae_blue}}{\thesection.}{0.6em}{}
\titleformat{\subsection}
  {\normalfont\normalsize\bfseries\sffamily\color{cae_blue!85!black}}{\thesubsection}{0.6em}{}
\titlespacing*{\section}{0pt}{14pt}{6pt}
\titlespacing*{\subsection}{0pt}{10pt}{4pt}

% Tiêu đề không đánh số ở đầu trang: \section* sau \clearpage nằm sát lề trên
% nên đè lên banner letterhead. Dùng \vspace* ngay TRƯỚC \section* không đủ,
% vì titlesec chèn khoảng cách riêng. Cách chắc chắn là đệm bằng \null.
\newcommand{\sectiondautrang}[1]{\section*{#1}}

% Dạng bài dùng \subsection*{} nên cần định dạng riêng, đậm và có gạch chân màu
\titleformat{\subsubsection}
  {\normalfont\normalsize\bfseries\sffamily\color{cae_blue!85!black}}{\thesubsubsection}{0.6em}{}

% ---------- LETTERHEAD (ảnh nền toàn trang) ----------
% Ảnh letterhead của CAE Shanghai: banner xanh trên/dưới kèm logo chìm giữa.
% Tỉ lệ ảnh 1414x2000 = 0.7070, khớp A4 (0.7071) nên chèn thẳng không méo.
% Đặt ở lớp nền dưới cùng để chữ không bị ảnh che.
% Dùng \AddToShipoutPictureBG (không có dấu *) để áp dụng cho MỌI trang.
\usepackage{eso-pic}
\AddToShipoutPictureBG{%
  \AtPageLowerLeft{%
    \includegraphics[width=\paperwidth,height=\paperheight]{images/letterhead.png}%
  }%
}

% ---------- ĐẦU / CHÂN TRANG ----------
% Letterhead đã có banner trên và dưới, nên bỏ dòng chạy để tránh chồng chữ.
% Số trang đặt trong vùng nội dung, sát trên banner dưới, lề phải.
\pagestyle{fancy}
\setlength{\headheight}{0pt}
\fancyhf{}
\renewcommand{\headrulewidth}{0pt}
\renewcommand{\footrulewidth}{0pt}
\fancyfoot[R]{\raisebox{0.9cm}{\small\color{cae_blue}\thepage}}

% ---------- THỤT LỀ ĐOẠN ----------
% Mặc định LaTeX thụt lề 15pt nhưng bỏ qua đoạn ngay sau tiêu đề mục, nên
% trong cùng một mục có đoạn thụt đoạn không, nhìn lệch hàng. Tài liệu này
% ngắt đoạn bằng khoảng trắng chứ không bằng thụt lề, nên đặt parindent = 0
% và cho parskip giãn ra để vẫn phân biệt được các đoạn.
\setlength{\parindent}{0pt}
\setlength{\parskip}{0.6em}

% ---------- MÔI TRƯỜNG NỘI DUNG ----------
% Trong tcolorbox, LaTeX đặt parindent = parskip = 0, nên các đoạn văn trong
% khung dính liền nhau. Đặt lại khoảng cách đoạn cho mọi khung.
\newcommand{\khoangcachdoan}{\setlength{\parskip}{0.5em}\setlength{\parindent}{0pt}}
\tcbset{before upper={\khoangcachdoan}}

\newtcolorbox{dinhnghia}[1][]{
  colback=box_bg, colframe=box_line, boxrule=0.6pt, arc=1pt,
  fonttitle=\bfseries, title={Định nghĩa #1}, breakable}
\newtcolorbox{dinhly}[1][]{
  colback=box_bg, colframe=cae_gold, boxrule=0.6pt, arc=1pt,
  fonttitle=\bfseries, title={Tính chất #1}, breakable}
\newtcolorbox{vidu}[1][]{
  colback=white, colframe=black!40, boxrule=0.4pt, arc=1pt,
  fonttitle=\bfseries, title={Ví dụ #1}, breakable}
\newtcolorbox{baitap}[1][]{
  colback=white, colframe=cae_blue!50, boxrule=0.4pt, arc=1pt,
  fonttitle=\bfseries, title={Câu #1}, breakable}

% Ô đáp án: viền xanh đậm hơn, dùng trong file đáp án
\newtcolorbox{dapan}[1][]{
  colback=box_bg, colframe=cae_blue, boxrule=0.6pt, arc=1pt,
  fonttitle=\bfseries, title={Câu #1}, breakable}

% ---------- TÊN CHỨNG MINH ----------
\renewcommand{\proofname}{\normalfont\bfseries Chứng minh}

% ---------- MÔN HỌC ----------
% Đổi được ở file .tex trước khi % =====================================================================
%  CAE SHANGHAI — GÓI LỆNH DÙNG CHUNG CHO TÀI LIỆU LUYỆN THI CSCA TOÁN
%  Biên dịch bằng XeLaTeX. Mọi file .tex trong buổi học \input file này.
% =====================================================================

% ---------- NGÔN NGỮ & PHÔNG CHỮ ----------
\usepackage{fontspec}
\usepackage{polyglossia}
\setmainlanguage{vietnamese}
\setotherlanguage{english}

% Chữ Latin (Anh + Việt): TeX Gyre Termes — bản clone Times New Roman,
% đi kèm TeX Live, hỗ trợ đầy đủ dấu tiếng Việt.
\setmainfont{TeX Gyre Termes}
\setsansfont{TeX Gyre Heros}

% Chữ Hán: ưu tiên Source Han Serif SC (思源宋体) — chuẩn học thuật Trung Quốc.
% Nếu máy chưa cài, tự động dùng Songti SC (宋体, có sẵn trên macOS).
% Cả hai đều thuộc họ 宋体 nên diện mạo tài liệu không đổi.
\IfFontExistsTF{Source Han Serif SC}{
  \newcommand{\cjkfamilyname}{Source Han Serif SC}
}{
  \IfFontExistsTF{Songti SC}{
    \newcommand{\cjkfamilyname}{Songti SC}
  }{
    \newcommand{\cjkfamilyname}{Noto Serif CJK SC}
  }
}

% xeCJK tự chuyển font khi gặp ký tự Hán, không cần bọc \cjk{} thủ công.
% Đây là điều kiện bắt buộc: font Latin (TeX Gyre Termes) không có glyph chữ Hán,
% nếu không khai báo xeCJK thì mọi chữ Hán viết trần sẽ bị nuốt khỏi bản in.
\usepackage{xeCJK}
\setCJKfamilyfont{caehan}{\cjkfamilyname}[Scale=0.95]
\setCJKmainfont{\cjkfamilyname}[Scale=0.95]
\setCJKsansfont{\cjkfamilyname}[Scale=0.95]
\setCJKmonofont{\cjkfamilyname}[Scale=0.95]

% \cjk{} ép dùng font Hán trong ngữ cảnh cần thiết (ví dụ trong tiêu đề in đậm
% hoặc trong bảng có định dạng riêng). Dùng \CJKfamily chứ không phải tên font,
% vì tên font chỉ là văn bản — sẽ in ra chữ thay vì đổi font.
\newcommand{\cjk}[1]{{\CJKfamily{caehan}#1}}

% ---------- TOÁN HỌC ----------
\usepackage{amsmath, amssymb, amsthm, mathtools}
\usepackage{unicode-math}
\setmathfont{TeX Gyre Termes Math}

% ---------- BỐ CỤC TRANG ----------
% Lề phải chừa vùng an toàn của letterhead:
%   banner trên chiếm 0–19mm  → top tối thiểu 2,2cm
%   banner dưới chiếm 281–297mm → bottom tối thiểu 2,2cm
% Lề trên 2,6cm (26mm) để tiêu đề đầu trang không đè banner (banner hết ở 19mm).
\usepackage[a4paper, top=2.6cm, bottom=2.4cm, left=2.5cm, right=2.0cm]{geometry}
\usepackage{fancyhdr}
\usepackage{titlesec}
\usepackage{enumitem}
\usepackage{graphicx}
\usepackage{booktabs, array, longtable, tabularx}
\usepackage{xcolor}
\usepackage[most]{tcolorbox}
\usepackage[colorlinks=true, linkcolor=cae_blue, citecolor=cae_blue, urlcolor=cae_blue]{hyperref}

% ---------- BẢNG MÀU THƯƠNG HIỆU ----------
\definecolor{cae_blue}{RGB}{0, 51, 102}
\definecolor{cae_gold}{RGB}{196, 149, 60}
\definecolor{box_bg}{RGB}{245, 247, 250}
\definecolor{box_line}{RGB}{0, 51, 102}

% ---------- TIÊU ĐỀ MỤC ----------
% Đệm đầu trang: banner letterhead chiếm 0–19mm, mà lề trên là 22mm.
% \section* đầu trang (sau \clearpage) sát lề nên phải chừa thêm, nếu không
% sẽ đè lên banner.
\titleformat{\section}
  {\normalfont\large\bfseries\sffamily\color{cae_blue}}{\thesection.}{0.6em}{}
\titleformat{\subsection}
  {\normalfont\normalsize\bfseries\sffamily\color{cae_blue!85!black}}{\thesubsection}{0.6em}{}
\titlespacing*{\section}{0pt}{14pt}{6pt}
\titlespacing*{\subsection}{0pt}{10pt}{4pt}

% Tiêu đề không đánh số ở đầu trang: \section* sau \clearpage nằm sát lề trên
% nên đè lên banner letterhead. Dùng \vspace* ngay TRƯỚC \section* không đủ,
% vì titlesec chèn khoảng cách riêng. Cách chắc chắn là đệm bằng \null.
\newcommand{\sectiondautrang}[1]{\section*{#1}}

% Dạng bài dùng \subsection*{} nên cần định dạng riêng, đậm và có gạch chân màu
\titleformat{\subsubsection}
  {\normalfont\normalsize\bfseries\sffamily\color{cae_blue!85!black}}{\thesubsubsection}{0.6em}{}

% ---------- LETTERHEAD (ảnh nền toàn trang) ----------
% Ảnh letterhead của CAE Shanghai: banner xanh trên/dưới kèm logo chìm giữa.
% Tỉ lệ ảnh 1414x2000 = 0.7070, khớp A4 (0.7071) nên chèn thẳng không méo.
% Đặt ở lớp nền dưới cùng để chữ không bị ảnh che.
% Dùng \AddToShipoutPictureBG (không có dấu *) để áp dụng cho MỌI trang.
\usepackage{eso-pic}
\AddToShipoutPictureBG{%
  \AtPageLowerLeft{%
    \includegraphics[width=\paperwidth,height=\paperheight]{images/letterhead.png}%
  }%
}

% ---------- ĐẦU / CHÂN TRANG ----------
% Letterhead đã có banner trên và dưới, nên bỏ dòng chạy để tránh chồng chữ.
% Số trang đặt trong vùng nội dung, sát trên banner dưới, lề phải.
\pagestyle{fancy}
\setlength{\headheight}{0pt}
\fancyhf{}
\renewcommand{\headrulewidth}{0pt}
\renewcommand{\footrulewidth}{0pt}
\fancyfoot[R]{\raisebox{0.9cm}{\small\color{cae_blue}\thepage}}

% ---------- THỤT LỀ ĐOẠN ----------
% Mặc định LaTeX thụt lề 15pt nhưng bỏ qua đoạn ngay sau tiêu đề mục, nên
% trong cùng một mục có đoạn thụt đoạn không, nhìn lệch hàng. Tài liệu này
% ngắt đoạn bằng khoảng trắng chứ không bằng thụt lề, nên đặt parindent = 0
% và cho parskip giãn ra để vẫn phân biệt được các đoạn.
\setlength{\parindent}{0pt}
\setlength{\parskip}{0.6em}

% ---------- MÔI TRƯỜNG NỘI DUNG ----------
% Trong tcolorbox, LaTeX đặt parindent = parskip = 0, nên các đoạn văn trong
% khung dính liền nhau. Đặt lại khoảng cách đoạn cho mọi khung.
\newcommand{\khoangcachdoan}{\setlength{\parskip}{0.5em}\setlength{\parindent}{0pt}}
\tcbset{before upper={\khoangcachdoan}}

\newtcolorbox{dinhnghia}[1][]{
  colback=box_bg, colframe=box_line, boxrule=0.6pt, arc=1pt,
  fonttitle=\bfseries, title={Định nghĩa #1}, breakable}
\newtcolorbox{dinhly}[1][]{
  colback=box_bg, colframe=cae_gold, boxrule=0.6pt, arc=1pt,
  fonttitle=\bfseries, title={Tính chất #1}, breakable}
\newtcolorbox{vidu}[1][]{
  colback=white, colframe=black!40, boxrule=0.4pt, arc=1pt,
  fonttitle=\bfseries, title={Ví dụ #1}, breakable}
\newtcolorbox{baitap}[1][]{
  colback=white, colframe=cae_blue!50, boxrule=0.4pt, arc=1pt,
  fonttitle=\bfseries, title={Câu #1}, breakable}

% Ô đáp án: viền xanh đậm hơn, dùng trong file đáp án
\newtcolorbox{dapan}[1][]{
  colback=box_bg, colframe=cae_blue, boxrule=0.6pt, arc=1pt,
  fonttitle=\bfseries, title={Câu #1}, breakable}

% ---------- TÊN CHỨNG MINH ----------
\renewcommand{\proofname}{\normalfont\bfseries Chứng minh}

% ---------- MÔN HỌC ----------
% Đổi được ở file .tex trước khi \input{cae-style} bằng \renewcommand{\monhoc}{...}
\providecommand{\monhoc}{Môn Toán}

% ---------- TRANG BÌA GỌN ----------
% #1 = loại tài liệu (TÀI LIỆU CHUẨN BỊ, TÀI LIỆU BUỔI HỌC, ...)
% #2 = số buổi
% #3 = tên buổi tiếng Việt
% #4 = tên buổi tiếng Hán
% #5 = module và tỉ trọng
\newcommand{\trangbia}[5]{%
  \begin{center}
    \vspace*{0.6cm}
    {\large\bfseries\color{cae_blue} CAE SHANGHAI}\\[0.15cm]
    {\footnotesize Tài liệu luyện thi CSCA · \monhoc}\\[0.5cm]
    {\color{cae_gold}\rule{0.55\linewidth}{0.8pt}}\\[0.35cm]
    {\small\bfseries\color{cae_blue!80!black} #1}\\[0.35cm]
    {\Large\bfseries\color{cae_blue} BUỔI #2 — #3}\\[0.2cm]
    {\large\color{cae_blue!85!black} \cjk{#4}}\\[0.25cm]
    \ifx\relax#5\relax\else{\small #5}\\[0.35cm]\fi
    {\color{cae_gold}\rule{0.55\linewidth}{0.8pt}}\\[0.3cm]
    {\footnotesize Biên soạn: Hoàng Dương, CAE Shanghai \quad·\quad \today}
  \end{center}
  \vspace{0.3cm}
}

% ---------- TRANG BÌA CHO ĐỀ KIỂM TRA ----------
% Khác \trangbia ở chỗ không có chữ "BUỔI N" — dùng cho đề thi, không phải bài giảng.
% #1 = loại tài liệu, #2 = tên kỳ thi (Việt), #3 = tên kỳ thi (Hán), #4 = thông số
\newcommand{\biade}[4]{%
  \begin{center}
    \vspace*{0.6cm}
    {\large\bfseries\color{cae_blue} CAE SHANGHAI}\\[0.15cm]
    {\footnotesize Tài liệu luyện thi CSCA · \monhoc}\\[0.5cm]
    {\color{cae_gold}\rule{0.55\linewidth}{0.8pt}}\\[0.35cm]
    {\small\bfseries\color{cae_blue!80!black} #1}\\[0.35cm]
    {\Large\bfseries\color{cae_blue} #2}\\[0.2cm]
    {\large\color{cae_blue!85!black} \cjk{#3}}\\[0.25cm]
    {\small #4}\\[0.35cm]
    {\color{cae_gold}\rule{0.55\linewidth}{0.8pt}}\\[0.3cm]
    {\footnotesize Biên soạn: Hoàng Dương, CAE Shanghai \quad·\quad \today}
  \end{center}
  \vspace{0.3cm}
}

% ---------- TIỆN ÍCH ----------
% Nhãn từ khóa in đậm ở đầu đoạn, dùng cho Từ vựng / Từ khóa nhận dạng / Chú ý
% Khoảng cách đoạn do khung tự đặt (xem \khoangcachdoan ở trên),
% ở đây chỉ cần tách đoạn và bỏ thụt đầu dòng.
\newcommand{\nhan}[1]{\par\noindent\textbf{#1.}\ \ignorespaces}

% Bảng gọn, không đánh số
\newcommand{\bang}[1]{%
  \begin{center}#1\end{center}}

% Đánh dấu ô trống cho học sinh điền. Dùng \dotfill để đường chấm giãn hết
% bề ngang ô: ô "Nghĩa" trong bảng từ vựng rộng hơn 200pt, để 3.2cm cố định
% thì phần còn lại trống trơn, nhìn như cột bị bóp.
\newcommand{\chotrong}{\dotfill}
 bằng \renewcommand{\monhoc}{...}
\providecommand{\monhoc}{Môn Toán}

% ---------- TRANG BÌA GỌN ----------
% #1 = loại tài liệu (TÀI LIỆU CHUẨN BỊ, TÀI LIỆU BUỔI HỌC, ...)
% #2 = số buổi
% #3 = tên buổi tiếng Việt
% #4 = tên buổi tiếng Hán
% #5 = module và tỉ trọng
\newcommand{\trangbia}[5]{%
  \begin{center}
    \vspace*{0.6cm}
    {\large\bfseries\color{cae_blue} CAE SHANGHAI}\\[0.15cm]
    {\footnotesize Tài liệu luyện thi CSCA · \monhoc}\\[0.5cm]
    {\color{cae_gold}\rule{0.55\linewidth}{0.8pt}}\\[0.35cm]
    {\small\bfseries\color{cae_blue!80!black} #1}\\[0.35cm]
    {\Large\bfseries\color{cae_blue} BUỔI #2 — #3}\\[0.2cm]
    {\large\color{cae_blue!85!black} \cjk{#4}}\\[0.25cm]
    \ifx\relax#5\relax\else{\small #5}\\[0.35cm]\fi
    {\color{cae_gold}\rule{0.55\linewidth}{0.8pt}}\\[0.3cm]
    {\footnotesize Biên soạn: Hoàng Dương, CAE Shanghai \quad·\quad \today}
  \end{center}
  \vspace{0.3cm}
}

% ---------- TRANG BÌA CHO ĐỀ KIỂM TRA ----------
% Khác \trangbia ở chỗ không có chữ "BUỔI N" — dùng cho đề thi, không phải bài giảng.
% #1 = loại tài liệu, #2 = tên kỳ thi (Việt), #3 = tên kỳ thi (Hán), #4 = thông số
\newcommand{\biade}[4]{%
  \begin{center}
    \vspace*{0.6cm}
    {\large\bfseries\color{cae_blue} CAE SHANGHAI}\\[0.15cm]
    {\footnotesize Tài liệu luyện thi CSCA · \monhoc}\\[0.5cm]
    {\color{cae_gold}\rule{0.55\linewidth}{0.8pt}}\\[0.35cm]
    {\small\bfseries\color{cae_blue!80!black} #1}\\[0.35cm]
    {\Large\bfseries\color{cae_blue} #2}\\[0.2cm]
    {\large\color{cae_blue!85!black} \cjk{#3}}\\[0.25cm]
    {\small #4}\\[0.35cm]
    {\color{cae_gold}\rule{0.55\linewidth}{0.8pt}}\\[0.3cm]
    {\footnotesize Biên soạn: Hoàng Dương, CAE Shanghai \quad·\quad \today}
  \end{center}
  \vspace{0.3cm}
}

% ---------- TIỆN ÍCH ----------
% Nhãn từ khóa in đậm ở đầu đoạn, dùng cho Từ vựng / Từ khóa nhận dạng / Chú ý
% Khoảng cách đoạn do khung tự đặt (xem \khoangcachdoan ở trên),
% ở đây chỉ cần tách đoạn và bỏ thụt đầu dòng.
\newcommand{\nhan}[1]{\par\noindent\textbf{#1.}\ \ignorespaces}

% Bảng gọn, không đánh số
\newcommand{\bang}[1]{%
  \begin{center}#1\end{center}}

% Đánh dấu ô trống cho học sinh điền. Dùng \dotfill để đường chấm giãn hết
% bề ngang ô: ô "Nghĩa" trong bảng từ vựng rộng hơn 200pt, để 3.2cm cố định
% thì phần còn lại trống trơn, nhìn như cột bị bóp.
\newcommand{\chotrong}{\dotfill}
 bằng \renewcommand{\monhoc}{...}
\providecommand{\monhoc}{Môn Toán}

% ---------- TRANG BÌA GỌN ----------
% #1 = loại tài liệu (TÀI LIỆU CHUẨN BỊ, TÀI LIỆU BUỔI HỌC, ...)
% #2 = số buổi
% #3 = tên buổi tiếng Việt
% #4 = tên buổi tiếng Hán
% #5 = module và tỉ trọng
\newcommand{\trangbia}[5]{%
  \begin{center}
    \vspace*{0.6cm}
    {\large\bfseries\color{cae_blue} CAE SHANGHAI}\\[0.15cm]
    {\footnotesize Tài liệu luyện thi CSCA · \monhoc}\\[0.5cm]
    {\color{cae_gold}\rule{0.55\linewidth}{0.8pt}}\\[0.35cm]
    {\small\bfseries\color{cae_blue!80!black} #1}\\[0.35cm]
    {\Large\bfseries\color{cae_blue} BUỔI #2 — #3}\\[0.2cm]
    {\large\color{cae_blue!85!black} \cjk{#4}}\\[0.25cm]
    \ifx\relax#5\relax\else{\small #5}\\[0.35cm]\fi
    {\color{cae_gold}\rule{0.55\linewidth}{0.8pt}}\\[0.3cm]
    {\footnotesize Biên soạn: Hoàng Dương, CAE Shanghai \quad·\quad \today}
  \end{center}
  \vspace{0.3cm}
}

% ---------- TRANG BÌA CHO ĐỀ KIỂM TRA ----------
% Khác \trangbia ở chỗ không có chữ "BUỔI N" — dùng cho đề thi, không phải bài giảng.
% #1 = loại tài liệu, #2 = tên kỳ thi (Việt), #3 = tên kỳ thi (Hán), #4 = thông số
\newcommand{\biade}[4]{%
  \begin{center}
    \vspace*{0.6cm}
    {\large\bfseries\color{cae_blue} CAE SHANGHAI}\\[0.15cm]
    {\footnotesize Tài liệu luyện thi CSCA · \monhoc}\\[0.5cm]
    {\color{cae_gold}\rule{0.55\linewidth}{0.8pt}}\\[0.35cm]
    {\small\bfseries\color{cae_blue!80!black} #1}\\[0.35cm]
    {\Large\bfseries\color{cae_blue} #2}\\[0.2cm]
    {\large\color{cae_blue!85!black} \cjk{#3}}\\[0.25cm]
    {\small #4}\\[0.35cm]
    {\color{cae_gold}\rule{0.55\linewidth}{0.8pt}}\\[0.3cm]
    {\footnotesize Biên soạn: Hoàng Dương, CAE Shanghai \quad·\quad \today}
  \end{center}
  \vspace{0.3cm}
}

% ---------- TIỆN ÍCH ----------
% Nhãn từ khóa in đậm ở đầu đoạn, dùng cho Từ vựng / Từ khóa nhận dạng / Chú ý
% Khoảng cách đoạn do khung tự đặt (xem \khoangcachdoan ở trên),
% ở đây chỉ cần tách đoạn và bỏ thụt đầu dòng.
\newcommand{\nhan}[1]{\par\noindent\textbf{#1.}\ \ignorespaces}

% Bảng gọn, không đánh số
\newcommand{\bang}[1]{%
  \begin{center}#1\end{center}}

% Đánh dấu ô trống cho học sinh điền. Dùng \dotfill để đường chấm giãn hết
% bề ngang ô: ô "Nghĩa" trong bảng từ vựng rộng hơn 200pt, để 3.2cm cố định
% thì phần còn lại trống trơn, nhìn như cột bị bóp.
\newcommand{\chotrong}{\dotfill}


\begin{document}

\trangbia{TÀI LIỆU BUỔI HỌC}{1}{Tập hợp \& Bất đẳng thức}{集合与不等式}{Module M1 \quad·\quad Tỉ trọng đề thi $\sim$10\% \quad·\quad khoảng 4--6 câu}

% =====================================================================
\section{Từ khóa nhận dạng đề}

Đọc được từ khóa là xác định được dạng bài mà không cần dịch cả đề. Bảng dưới đây gom những từ khóa quyết định của chuyên đề này.

\bang{
\begin{tabular}{@{}lll@{}}
\toprule
\textbf{Từ khóa} & \textbf{Nghĩa} & \textbf{Dạng bài tương ứng} \\
\midrule
交集 & giao & Tìm phần chung của hai tập hợp \\
并集 & hợp & Gộp tất cả phần tử của hai tập \\
补集 & phần bù & Lấy phần của tập vũ trụ không thuộc tập đã cho \\
全集 & tập vũ trụ & Tập nền để lấy phần bù \\
子集 & tập con & Đếm số tập con: $2^n$ \\
真子集 & tập con thực sự & Đếm số tập con thực sự: $2^n - 1$ \\
元素个数 & số phần tử & Đếm phần tử của tập hợp \\
解集 & tập nghiệm & Giải bất phương trình, viết kết quả dạng tập hợp \\
取值范围 & khoảng giá trị & Bài toán chứa tham số \\
恒成立 & đúng với mọi giá trị của biến & Tìm khoảng tham số để biểu thức luôn đúng \\
有解 & có nghiệm & Tồn tại ít nhất một giá trị thỏa mãn \\
至少 & ít nhất & Đếm phần bù hoặc dùng biến cố đối \\
\bottomrule
\end{tabular}}

\nhan{Chú ý} Hai từ khóa 恒成立 (đúng với mọi $x$) và 有解 (có nghiệm) ngược nhau về cách xử lý. Gặp 恒成立 thì dùng điều kiện $\Delta$; gặp 有解 thì tìm giá trị nhỏ nhất hoặc lớn nhất của biểu thức. Đọc nhầm hai từ này là mất điểm chắc chắn.

% =====================================================================
\section{Kiến thức trọng tâm}

\subsection{Tập hợp và quan hệ}

\begin{dinhnghia}[1 — Tập hợp con]
Tập hợp $A$ được gọi là \emph{tập hợp con} của $B$ (kí hiệu $A \subseteq B$) khi và chỉ khi mọi phần tử của $A$ đều là phần tử của $B$.
\end{dinhnghia}

\begin{dinhnghia}[2 — Tập hợp bằng nhau]
Hai tập hợp $A$ và $B$ bằng nhau khi và chỉ khi $A \subseteq B$ và $B \subseteq A$.
\end{dinhnghia}

\begin{dinhly}[1 — Các tập hợp số]
\[
\mathbb{N} \subseteq \mathbb{Z} \subseteq \mathbb{Q} \subseteq \mathbb{R}
\]
\end{dinhly}

\begin{dinhly}[2 — Số tập con]
Tập hợp có $n$ phần tử có $2^n$ tập con, trong đó có $2^n - 1$ tập con thực sự.
\end{dinhly}

\nhan{Chú ý} Phân biệt $A \subseteq B$ (cho phép $A = B$) và $A \subsetneq B$ (không cho phép $A = B$). Đây là lỗi thường gặp khi đề hỏi 真子集.

\nhan{Chú ý} Tập rỗng $\varnothing$ là tập con của mọi tập hợp, và là tập con thực sự của mọi tập hợp khác rỗng. Đề rất hay khai thác điểm này.

\subsection{Các phép toán tập hợp}

\begin{dinhnghia}[3 — Giao]
$A \cap B = \{x \mid x \in A \ \text{và} \ x \in B\}$.
\end{dinhnghia}

\begin{dinhnghia}[4 — Hợp]
$A \cup B = \{x \mid x \in A \ \text{hoặc} \ x \in B\}$.
\end{dinhnghia}

\begin{dinhnghia}[5 — Phần bù]
$\complement_U A = \{x \mid x \in U \ \text{và} \ x \notin A\}$.
\end{dinhnghia}

\begin{dinhly}[3 — Luật De Morgan]
\[
\complement_U (A \cup B) = \complement_U A \cap \complement_U B
\]
\[
\complement_U (A \cap B) = \complement_U A \cup \complement_U B
\]
\end{dinhly}

\nhan{Chú ý} Luật De Morgan là công cụ rút gọn mạnh nhất khi đề cho 补集 của một biểu thức phức tạp. Thay vì tính $A \cup B$ rồi lấy phần bù, có thể lấy phần bù từng tập rồi giao lại.

\begin{dinhly}[4 — Quan hệ tương đương với tập con]
\[
A \cap B = B \iff B \subseteq A, \qquad A \cup B = A \iff B \subseteq A
\]
\end{dinhly}

\nhan{Chú ý} Hai quan hệ này biến bài toán ``tìm tham số để $A \cap B = B$'' thành bài toán tập con quen thuộc. Nhận ra phép biến đổi này tiết kiệm rất nhiều thời gian.

\subsection{Bất đẳng thức và bất phương trình}

\begin{dinhnghia}[6 — Bất phương trình bậc hai]
Bất phương trình dạng $ax^2 + bx + c > 0$ (hoặc $<, \ge, \le$) với $a \neq 0$.
\end{dinhnghia}

\begin{dinhly}[5 — Quy tắc xét dấu tam thức bậc hai]
Xét tam thức $f(x) = ax^2 + bx + c$ với $a > 0$, gọi $\Delta = b^2 - 4ac$.

\bang{
\begin{tabular}{@{}lcc@{}}
\toprule
\textbf{Trường hợp} & $f(x) > 0$ & $f(x) < 0$ \\
\midrule
$\Delta > 0$, hai nghiệm $x_1 < x_2$ & $x < x_1$ hoặc $x > x_2$ & $x_1 < x < x_2$ \\
$\Delta = 0$, nghiệm kép $x_0$ & $x \neq x_0$ & vô nghiệm \\
$\Delta < 0$ & mọi $x \in \mathbb{R}$ & vô nghiệm \\
\bottomrule
\end{tabular}}
\end{dinhly}

\nhan{Chú ý} Quy tắc ghi nhớ: \textbf{``trong trái, ngoài cùng''} — với $a > 0$, biểu thức cùng dấu với $a$ ở ngoài khoảng hai nghiệm, trái dấu ở trong khoảng hai nghiệm. Nếu $a < 0$, nhân cả hai vế với $-1$ và \textbf{đổi chiều bất đẳng thức} rồi áp dụng bảng trên.

\begin{dinhly}[6 — Bất phương trình chứa dấu giá trị tuyệt đối]
Với $a > 0$:
\[
|f(x)| < a \iff -a < f(x) < a
\]
\[
|f(x)| > a \iff f(x) < -a \ \text{hoặc} \ f(x) > a
\]
\end{dinhly}

\begin{dinhly}[7 — Bất phương trình tích và thương]
Đưa về dạng $f(x) \cdot g(x) > 0$ hoặc $\dfrac{f(x)}{g(x)} > 0$, xét dấu từng nhân tử trên trục số rồi lấy miền thỏa mãn.
\end{dinhly}

\nhan{Chú ý} Với bất phương trình thương, \textbf{nghiệm của mẫu luôn bị loại}, kể cả khi đề dùng dấu $\ge$ hay $\le$. Đây là lỗi mất điểm phổ biến nhất.

\begin{dinhly}[8 — Bất phương trình bậc hai đúng với mọi $x$]
$ax^2 + bx + c > 0$ đúng với mọi $x \in \mathbb{R}$ khi và chỉ khi $a > 0$ và $\Delta < 0$.
\end{dinhly}

\nhan{Chú ý} Nếu đề không khẳng định $a \neq 0$, \textbf{phải xét riêng trường hợp $a = 0$}. Đây chính là bẫy của dạng bài này.

\begin{dinhly}[9 — Bất đẳng thức Cauchy (AM--GM)]
Với $a, b \ge 0$:
\[
a + b \ge 2\sqrt{ab}
\]
Dấu bằng xảy ra khi $a = b$.
\end{dinhly}

\nhan{Hệ quả} Với $a, b > 0$: $\dfrac{1}{a} + \dfrac{1}{b} \ge \dfrac{4}{a+b}$.

\nhan{Hệ quả} Với $a^2 + b^2 = R^2$: giá trị lớn nhất của $pa + qb$ là $R\sqrt{p^2 + q^2}$.

\clearpage
% =====================================================================
\section{Các dạng bài}

\subsection*{Dạng 1 — Phép toán trên tập hợp cho dưới dạng mô tả tính chất}

Cho hai tập hợp viết dưới dạng $\{x \mid \text{điều kiện}\}$, yêu cầu tìm giao, hợp hoặc phần bù.

\nhan{Cách nhận dạng} Đề cho từ hai tập hợp trở lên và yêu cầu tìm tập hợp mới. Các tập hợp được mô tả bằng bất đẳng thức hoặc bất phương trình.

\nhan{Cách làm} Giải từng tập hợp về dạng khoảng trên trục số, vẽ trục số, rồi lấy phần chung (giao) hoặc phần gộp (hợp).

\begin{vidu}[1]
Cho $A = \{x \mid -1 \le x < 10\}$ và $B = \{x \mid x \ge 2\}$. Tìm $A \cap B$.

\textbf{Giải.} Biểu diễn $A = [-1, 10)$ và $B = [2, +\infty)$ trên trục số. Phần chung là từ $2$ đến dưới $10$, bao gồm cả $2$.
\[
A \cap B = [2, 10) = \{x \mid 2 \le x < 10\}
\]
\end{vidu}

\subsection*{Dạng 2 — Đếm số phần tử và số tập con}

\nhan{Cách nhận dạng} Đề hỏi 元素个数 (số phần tử), 子集个数 (số tập con), 真子集个数 (số tập con thực sự). Thường kèm điều kiện như $x \in \mathbb{Z}$ để tập hợp là hữu hạn.

\nhan{Cách làm} Đếm số phần tử $n$ của tập hợp, sau đó áp dụng $2^n$ (tập con) hoặc $2^n - 1$ (tập con thực sự). Với tập hợp cho dạng khoảng, đếm số nguyên trong khoảng đó.

\begin{vidu}[2]
Cho tập hợp $A = \{x \in \mathbb{Z} \mid |x| < 2\}$. Tìm số tập con của $A$.

\textbf{Giải.} Điều kiện $|x| < 2$ tương đương $-2 < x < 2$. Vì $x \in \mathbb{Z}$ nên $A = \{-1, 0, 1\}$, có $n = 3$ phần tử.

Số tập con là $2^3 = 8$.
\end{vidu}

\subsection*{Dạng 3 — Kết hợp tập hợp với bất phương trình}

\nhan{Cách nhận dạng} Mỗi tập hợp được cho bởi một bất phương trình (bậc hai, chứa dấu giá trị tuyệt đối, hoặc phân thức) chứ không phải khoảng cho sẵn. Đây là dạng phổ biến nhất trong đề thi.

\nhan{Cách làm} Giải từng bất phương trình để đưa tập hợp về dạng khoảng, rồi thực hiện phép toán tập hợp. Thứ tự bắt buộc: \textbf{giải bất phương trình trước, vẽ trục số sau}.

\begin{vidu}[3]
Cho $A = \{x \mid x^2 - 2x - 3 > 0\}$ và $B = \{x \mid x^2 - 4 < 0\}$. Tìm $A \cap B$.

\textbf{Giải.} Giải $x^2 - 2x - 3 > 0$: phương trình $x^2 - 2x - 3 = 0$ có hai nghiệm $x = -1$ và $x = 3$. Vì $a = 1 > 0$ nên nghiệm của bất phương trình là ngoài khoảng hai nghiệm:
\[
A = (-\infty, -1) \cup (3, +\infty)
\]
Giải $x^2 - 4 < 0$: hai nghiệm $x = \pm 2$, vì $a = 1 > 0$ nên nghiệm là trong khoảng hai nghiệm:
\[
B = (-2, 2)
\]
Giao hai tập hợp: phần chung của $(-\infty, -1) \cup (3, +\infty)$ với $(-2, 2)$ là $(-2, -1)$.
\[
A \cap B = (-2, -1)
\]
\end{vidu}

\subsection*{Dạng 4 — Bất phương trình chứa dấu giá trị tuyệt đối}

\nhan{Cách nhận dạng} Đề có dạng $|f(x)| < a$, $|f(x)| > a$, hoặc biểu thức chứa hai dấu giá trị tuyệt đối như $|x - 2| + |x + 1| < 5$.

\nhan{Cách làm} Với một dấu, dùng Tính chất 6 chuyển về bất phương trình không chứa dấu. Với hai dấu, chia khoảng theo các điểm làm biểu thức trong dấu đổi dấu (điểm 0 của từng biểu thức), giải trên từng khoảng rồi hợp nghiệm.

\begin{vidu}[4]
Giải bất phương trình $|2x - 1| \le 3$.

\textbf{Giải.} Áp dụng $|f(x)| \le a \iff -a \le f(x) \le a$ với $a = 3$:
\[
-3 \le 2x - 1 \le 3
\]
Cộng $1$ vào cả ba vế: $-2 \le 2x \le 4$. Chia cho $2$: $-1 \le x \le 2$.

Tập nghiệm là $[-1, 2]$.
\end{vidu}

\subsection*{Dạng 5 — Bất phương trình tích và bất phương trình thương}

\nhan{Cách nhận dạng} Đề có dạng tích hai biểu thức như $(x^2 - 1)(x - 2) < 0$, hoặc phân thức như $\dfrac{x - 4}{4x^2 - 1} < 0$, $\dfrac{x - 3}{x + 2} \ge 2$.

\nhan{Cách làm} Đưa về dạng $f(x) \cdot g(x)$ hoặc $\dfrac{f(x)}{g(x)}$ với một vế bằng $0$. Tìm tất cả nghiệm của tử và mẫu, sắp xếp trên trục số, xét dấu trên từng khoảng, chọn khoảng thỏa mãn. Nghiệm của mẫu luôn bị loại.

\begin{vidu}[5]
Giải bất phương trình $\dfrac{x - 4}{4x^2 - 1} < 0$.

\textbf{Giải.} Mẫu $4x^2 - 1 = (2x-1)(2x+1)$ có hai nghiệm $x = \dfrac{1}{2}$ và $x = -\dfrac{1}{2}$; tử có nghiệm $x = 4$. Ba điểm chia trục số thành bốn khoảng.

\bang{
\begin{tabular}{@{}lccc@{}}
\toprule
\textbf{Khoảng} & \textbf{Tử} $x-4$ & \textbf{Mẫu} $4x^2-1$ & \textbf{Thương} \\
\midrule
$x < -\dfrac{1}{2}$ & $-$ & $+$ & $-$ \checkmark \\
$-\dfrac{1}{2} < x < \dfrac{1}{2}$ & $-$ & $-$ & $+$ \\
$\dfrac{1}{2} < x < 4$ & $-$ & $+$ & $-$ \checkmark \\
$x > 4$ & $+$ & $+$ & $+$ \\
\bottomrule
\end{tabular}}

Tập nghiệm là $\left(-\infty, -\dfrac{1}{2}\right) \cup \left(\dfrac{1}{2}, 4\right)$.
\end{vidu}

\subsection*{Dạng 6 — Bài toán chứa tham số và điều kiện tập con}

\nhan{Cách nhận dạng} Đề có tham số $a$, $m$, $k$ và yêu cầu tìm giá trị để $A \subseteq B$, hoặc $A \cap B = B$, hoặc $A \cup B = A$.

\nhan{Cách làm} Dùng Tính chất 4 để đưa về điều kiện tập con, sau đó so sánh hai đầu mút của khoảng. \textbf{Luôn xét trường hợp tập hợp rỗng trước} — nhiều bài toán có đáp án đúng nằm ở trường hợp này.

\begin{vidu}[6]
Cho $A = \{-1, \tfrac{1}{2}\}$ và $B = \{x \mid ax + 1 = 0\}$. Tìm $a$ để $A \cap B = B$.

\textbf{Giải.} Điều kiện $A \cap B = B$ tương đương $B \subseteq A$.

Xét $a = 0$: phương trình $ax + 1 = 0$ trở thành $1 = 0$, vô nghiệm nên $B = \varnothing$. Tập rỗng là tập con của mọi tập hợp, vậy $a = 0$ thỏa mãn.

Xét $a \neq 0$: phương trình có nghiệm duy nhất $x = -\dfrac{1}{a}$. Để $B \subseteq A$ thì $-\dfrac{1}{a}$ phải bằng $-1$ hoặc $\dfrac{1}{2}$.

Với $-\dfrac{1}{a} = -1$ được $a = 1$. Với $-\dfrac{1}{a} = \dfrac{1}{2}$ được $a = -2$.

Vậy $a \in \{-2, 0, 1\}$.
\end{vidu}

\subsection*{Dạng 7 — Bất đẳng thức đúng với mọi giá trị của biến}

\nhan{Cách nhận dạng} Đề có cụm 恒成立 (đúng với mọi $x$) hoặc 对于一切实数 $x$ (với mọi số thực $x$), yêu cầu tìm khoảng giá trị của tham số.

\nhan{Cách làm} Với $ax^2 + bx + c > 0$ đúng với mọi $x$, điều kiện là $a > 0$ và $\Delta < 0$. Với dấu $\ge$, điều kiện là $a > 0$ và $\Delta \le 0$. Nếu đề không khẳng định $a \neq 0$, phải xét riêng trường hợp $a = 0$.

\begin{vidu}[7]
Tìm $m$ để $x^2 + (m - 1)x + 1 > 0$ đúng với mọi $x \in \mathbb{R}$.

\textbf{Giải.} Hệ số $a = 1 > 0$ nên chỉ cần điều kiện $\Delta < 0$:
\[
\Delta = (m-1)^2 - 4 \cdot 1 \cdot 1 < 0 \iff (m-1)^2 < 4
\]
\[
\iff -2 < m - 1 < 2 \iff -1 < m < 3
\]
Vậy $m \in (-1, 3)$.
\end{vidu}

\clearpage
% =====================================================================
\section*{Phụ lục — Bối cảnh kỳ thi}
\addcontentsline{toc}{section}{Phụ lục — Bối cảnh kỳ thi}

\subsection*{Cấu trúc đề thi CSCA môn Toán}

\bang{
\begin{tabular}{@{}ll@{}}
\toprule
\textbf{Đặc điểm} & \textbf{Thông số} \\
\midrule
Số câu & 48 câu trắc nghiệm \\
Thời gian & 60 phút \\
Thang điểm & 100 điểm \\
Ngôn ngữ & Tiếng Trung hoặc tiếng Anh (thí sinh chọn) \\
Máy tính & Không được dùng \\
\bottomrule
\end{tabular}}

Thời gian trung bình mỗi câu là \textbf{1 phút 15 giây}. Mục tiêu của khóa học không phải giải được bài khó, mà là \textbf{nhận ra dạng bài trong vài giây và chọn đúng đáp án trong khoảng 40 giây}.

\subsection*{Trọng số các nhóm nội dung}

\bang{
\begin{tabular}{@{}lcc@{}}
\toprule
\textbf{Nhóm nội dung} & \textbf{Số câu} & \textbf{Tỉ trọng} \\
\midrule
Lượng giác & 8--11 & $\sim$20\% \\
Hình học giải tích & 8--10 & $\sim$19\% \\
Dãy số & 5--8 & $\sim$13\% \\
Hàm số & 5--7 & $\sim$13\% \\
\textbf{Tập hợp \& bất đẳng thức} & \textbf{4--6} & \textbf{$\sim$10\%} \\
Mũ \& logarit & 3--5 & $\sim$8\% \\
Vector \& số phức & 2--4 & $\sim$6\% \\
Xác suất \& thống kê & 1--2 & $\sim$3\% \\
\bottomrule
\end{tabular}}

Chuyên đề hôm nay chiếm khoảng 10\% đề thi — nhóm nội dung nhỏ, nhưng là \textbf{nền tảng của mọi nhóm khác}. Không đọc được ký hiệu tập hợp thì không giải được bất phương trình, tìm tập xác định hàm số, hay xét dấu biểu thức.

\vfill
\begin{center}
{\footnotesize\color{cae_blue} Tài liệu nội bộ — CAE SHANGHAI}
\end{center}

\end{document}
